% Three routes onto that bus, and what each one costs.
\begin{tikzpicture}[font=\sffamily\scriptsize, x=1cm, y=1cm,
  tinyblk/.style={blk, text width=17mm, minimum height=6mm, font=\sffamily\tiny,
                  inner sep=2pt}]

\node[chlabel, anchor=west] at (-6.8,3.4) {three routes onto that bus, none of them taken here};

% ================================================================ route 1
\node[layerlabel] at (-6.8,2.9) {1. keep this part, add a controller beside it};
\node[tinyblk, fill=hwcol] (a1) at (-5.6,2.0) {this board};
\node[tinyblk, fill=white] (a2) at (-3.4,2.0) {a controller chip};
\node[tinyblk, fill=black!6] (a3) at (-1.2,2.0) {memory, PHYs, magnetics};
\draw[flow] (a1) -- (a2);
\draw[busstub] (a2) -- (a3);
\node[chlabel, anchor=west, text width=44mm] at (0.2,2.0)
  {The most instructive route: fifteen chapters of this volume survive
   unchanged. \textbf{It needs soldering, which this bench cannot do.}};

% ================================================================ route 2
\node[layerlabel] at (-6.8,1.1) {2. replace this part with one that has a controller inside};
\node[tinyblk, fill=white] (b1) at (-5.6,0.2) {one chip: bus and core};
\node[tinyblk, fill=black!6] (b2) at (-3.4,0.2) {PHYs and magnetics};
\draw[busstub] (b1) -- (b2);
\node[chlabel, anchor=west, text width=44mm] at (0.2,0.2)
  {The cleanest route, and a different board with a different core, so the board
   support package, the clock tree and the memory map of chapters 1 to 4 all
   move.};

% ================================================================ route 3
\node[layerlabel] at (-6.8,-0.7) {3. a part that implements the protocol in programmable logic};
\node[tinyblk, fill=white] (c1) at (-5.6,-1.6) {programmable unit};
\node[tinyblk, fill=black!6] (c2) at (-3.4,-1.6) {a vendor's certified stack};
\draw[flow] (c1) -- (c2);
\node[chlabel, anchor=west, text width=44mm] at (0.2,-1.6)
  {Flexible, and switchable between protocols. \textbf{It ties the design to one
   vendor's stack}, which is a commercial decision rather than a technical one.};

% ================================================================ the one worth naming
\node[regcell, fill=usrcol, minimum width=114mm, text width=112mm, inner sep=3pt,
      anchor=north west, align=left] at (-6.8,-2.6)
  {\textbf{One controller in route 1 is built for exactly this problem}, and its
   own feature list says so with no interpretation needed: three channel pulse
   width modulation, step and direction control, incremental and Hall encoder
   interfaces, and an emergency stop input, with both physical layer devices
   integrated, in an eighty pin package. \textbf{That is a joint node's entire
   peripheral set, inside the bus chip.} Chapters 5, 7 and 18 would each have
   less work to do, and the node would have less of itself left.};

\node[umlnote, text width=114mm, anchor=north west] at (-6.8,-5.1)
  {No route is taken. The figure prices three purchases rather than describing a
   connection, because the chapter's conclusion is not to buy any of them yet,
   and a costed refusal is a more useful thing to leave in a repository than a
   half-finished attempt.};
\end{tikzpicture}
